\documentclass[11pt]{article}

\usepackage[utf8]{inputenc}
\usepackage[T1]{fontenc}
\usepackage{lmodern}
\usepackage{geometry}
\usepackage{amsmath,amssymb,amsfonts,amsthm,mathtools}
\usepackage{bm}
\usepackage{enumitem}
\usepackage{microtype}
\usepackage{hyperref}
\usepackage{titlesec}
\usepackage{booktabs}
\usepackage{array}
\usepackage{setspace}
\usepackage{mathrsfs}
\usepackage{physics}

\geometry{margin=1in}
\hypersetup{colorlinks=true, linkcolor=black, urlcolor=blue, citecolor=black}
\setlist[enumerate]{topsep=0.4em,itemsep=0.35em}
\setlist[itemize]{topsep=0.3em,itemsep=0.25em}
\titleformat{\section}{\large\bfseries}{\thesection.}{0.5em}{}
\titleformat{\subsection}{\normalsize\bfseries}{\thesubsection.}{0.5em}{}
\setstretch{1.06}

\newcommand{\Aether}{\AE{}ther}
\newcommand{\Aflow}{\AE{}ther-flow}
\newcommand{\TheoryName}{\Aflow{} Interpretation of Relativity}
\newcommand{\TheoryProgram}{\Aether{} / \Aflow{} framework}

\newcommand{\CompletedMetric}{g^{\sharp}}

\newcommand{\Car}{\mathrm{car}}
\newcommand{\Cur}{\mathrm{cur}}
\newcommand{\Hom}{\mathrm{hom}}

\newcommand{\ForSpace}[1]{\mathcal I_{#1}^{\mathrm{for}}}
\newcommand{\PrimShell}{\mathcal S_{\mathrm{prim}}}

\newcommand{\Reduction}[1]{\mathcal R_{#1}}
\newcommand{\CurrentCarrier}[1]{\mathfrak C_{#1}^{\mathrm{cur}}}
\newcommand{\EqCarrier}[1]{\mathfrak C_{#1}^{\mathrm{eq}}}
\newcommand{\CurrentExtract}[1]{\mathscr E_{#1}^{\mathrm{cur}}}
\newcommand{\EqExtract}[1]{\mathscr E_{#1}^{\mathrm{eq}}}
\newcommand{\PrimCurrent}[1]{\mathcal J_{#1}^{\mathrm{prim}}}
\newcommand{\PrimTheta}[1]{\Theta_{#1}^{\mathrm{prim}}}

\newcommand{\MetricOut}{\mathscr G_{\mathrm{out}}}
\newcommand{\MatterOut}{\mathscr T_{\mathrm{out}}}

\newcommand{\VisibleAffine}[1]{\widehat X_{#1}}
\newcommand{\DataSpace}[1]{\mathcal Y_{#1}^{\mathrm{obs}}}
\newcommand{\AmbientTarget}[1]{E_{#1}^{\mathrm{amb}}}

\newcommand{\GapSpace}[1]{H_{#1}}
\newcommand{\GapBody}[1]{D_{#1}}
\newcommand{\HiddenOperator}[1]{K_{#1}^{\mathrm{hid}}}

\newcommand{\GenSpace}[1]{\mathcal G_{#1}^{\mathrm{sel}}}
\newcommand{\GenBody}[1]{\mathcal B_{#1}^{\mathrm{sel}}}
\newcommand{\GenData}[1]{\Lambda_{#1}^{\mathrm{sel,dat}}}
\newcommand{\GenShell}[1]{\Lambda_{#1}^{\mathrm{sel,sh}}}
\newcommand{\HiddenSel}[1]{\Sigma_{#1}^{\mathrm{hid}}}
\newcommand{\ZeroHiddenSlice}[1]{\mathcal Z_{#1}^{0}}
\newcommand{\ZeroHiddenVec}[1]{V_{#1}^{0}}
\newcommand{\GenSym}[1]{\mathbb S_{#1}}
\newcommand{\GenTime}[1]{\vartheta_{#1}^{\mathrm{sel}}}
\newcommand{\GenEmbed}[1]{\zeta_{#1}^{\mathrm{sel}}}
\newcommand{\Kernel}[1]{K_{#1}^{0}}
\newcommand{\StateComp}[1]{P_{#1}^{\mathrm{co}}}

\newcommand{\SpecSpace}[1]{\mathcal U_{#1}^{\mathrm{sp}}}
\newcommand{\SpecBody}[1]{\mathcal B_{#1}^{\mathrm{sp}}}
\newcommand{\SpecData}[1]{\Lambda_{#1}^{\mathrm{sp,dat}}}
\newcommand{\SpecShell}[1]{\Lambda_{#1}^{\mathrm{sp,sh}}}
\newcommand{\SpecOp}[1]{\mathcal L_{#1}^{\mathrm{sp}}}
\newcommand{\SpecZero}[1]{V_{#1}^{\mathrm{sp},0}}
\newcommand{\SpecHidden}[1]{H_{#1}^{\mathrm{sp}}}
\newcommand{\SpecHiddenBody}[1]{D_{#1}^{\mathrm{sp}}}
\newcommand{\SpecProj}[1]{\Pi_{#1}^{\mathrm{sp,hid}}}
\newcommand{\SpecSym}[1]{\mathbb U_{#1}}
\newcommand{\SpecTime}[1]{\vartheta_{#1}^{\mathrm{sp}}}
\newcommand{\SpecEmbed}[1]{\zeta_{#1}^{\mathrm{sp}}}

\newtheorem{definition}{Definition}
\newtheorem{lemma}{Lemma}
\newtheorem{proposition}{Proposition}
\newtheorem{remark}{Remark}
\newtheorem{corollary}{Corollary}
\newtheorem{theorem}{Theorem}

\title{\textbf{The \TheoryName{}}\\[0.4em]
\large Primitive-Reservoir Observer-Reduction\\
\large Zero-Hidden Selector Generator Law\\
\large Necessity / Uniqueness from Zero-Hidden Spectral Generator Law}
\author{}
\date{}

\begin{document}

\maketitle

\begin{abstract}
The current benchmark-facing theorem on the active primitive-reservoir observer-reduction line already sharpened the burden once more: the completion-balance ambient dynamics package is no longer merely available, because it is derived canonically from a zero-hidden selector generator law and is necessary there up to the obvious selector-splitting identifications. The remaining honest task is therefore no longer to justify the ambient-balance layer. It is to justify why the selector law itself should hold.

The present note takes the smallest honest step now available below that selector layer without claiming final microphysical closure. Rather than postulate both a hidden selector and a separate hidden-sector operator, it collapses them into a single deeper zero-hidden spectral generator law. For each sector this smaller package consists of a compact convex spectral state body, affine observer-data and shell-response readouts, one positive self-adjoint spectral generator with an exact zero band and a positively gapped hidden band, exact body splitting over those spectral directions, symmetry and time-reversal actions preserving the spectral law, an exact zero-response shell realization on the zero band, fixed-data shell solvability on that same locus, and a spectral coercivity inequality on data-invisible variations modulo the inert zero-hidden kernel.

From that spectral law one derives canonically every constituent of the selector package that now requires justification: the compact convex selector-state body, the affine observer-data and shell-response readouts, the self-adjoint idempotent hidden selector, the selector-zero locus and hidden selector image, the selector-zero / hidden body splitting, the hidden factor and positive hidden-sector operator, the symmetry and time-reversal structure, the exact zero-response shell realization on the selector-zero locus, fixed-data shell solvability on that locus, the inert zero-hidden kernel and projector, and the selector-law coercivity inequality. Within this spectral class those selector constituents are not merely available. They are necessary and unique up to the obvious orthogonal identification preserving the zero/hidden spectral splitting and the recorded data and shell coordinates. The benchmark package remains fixed throughout: the one operative metric is still exactly \(\CompletedMetric\), the ordinary matter route is unchanged, there is no second operative metric, and the low-energy matter law is not enlarged.
\end{abstract}

\section{Introduction}

The active derivational program of \TheoryName{} has again narrowed what counts as an honest next step. The primitive-observer-reduction datum theorem, the defect-fiber factorization theorem, the one-metric output theorem, the chain-forcing theorem, the selector theorem, the stationary-construction theorem, the explicit-package theorem, the primitive-normal-form theorem, the ambient-shell-law theorem, the combined-response theorem, the graph-collapse theorem, the completion-core theorem, and the later ambient-dynamics theorem each discharged what was honestly open at its own layer. The latest result matters because it removes another tempting stopping point. Once the completion-balance ambient dynamics package is itself a canonical reduction of a deeper selector law, it is no longer enough to restate that ambient layer under a different presentation.

The remaining honest burden is therefore one layer deeper again. The task is now to explain why the zero-hidden selector generator law should exist, rather than letting the selector-state body, the data and shell readouts, the hidden selector, the selector-zero / hidden splitting, the hidden-sector operator, the shell realization, the solvability statement, and the coercive selector inequality stand as the new deepest disciplined assumptions. That is the burden addressed here.

The present note makes the smallest honest collapse that can presently be stated on the active line. It does not claim a final microscopic substrate law. Its gain is narrower and more upstream. It replaces the pair consisting of a hidden selector and a hidden-sector operator by a single positive self-adjoint spectral generator whose exact zero band and positively gapped hidden band canonically generate the selector package. In that precise sense, the selector law itself ceases to be an independent stopping point.

\paragraph{Framework and claim boundary.}
\TheoryProgram{} denotes the broader \Aether{} / \Aflow{} framework built on the claims that the \Aether{} is the underlying four-dimensional substrate of reality, the \Aflow{} is its intrinsic ordered motion, observed three-dimensional space is the local experiential slice of that deeper substrate, S-time is the experienced order of change arising from matter, light, and the \Aflow{}, and observed expansion is the three-dimensional appearance of deeper four-dimensional ordered motion. Within that framework, \TheoryName{} denotes the exact-closure theory in which the effective gravitational dynamics are adopted to be exactly Einsteinian with universal matter coupling. In the active sequence, \emph{adoption} means use of that established relativistic dynamics without claiming substrate derivation, while \emph{derivation} is reserved for a first-principles recovery from explicit substrate variables.

The present note stays strictly inside that boundary. It does not claim that final substrate microphysics has already been found. Its narrower theorem is only that, on the active primitive-observer-reduction line, the zero-hidden selector generator law is itself forced by a smaller zero-hidden spectral generator law, and is unique there up to the obvious orthogonal identification preserving the zero/hidden spectral splitting.

\section{Fixed Primitive Continuation Data}

\begin{definition}[Recorded primitive continuation data]
\label{def:fixed}
Fix the following continuation data from the active primitive-reservoir observer-reduction line.
\begin{enumerate}
    \item For each sector label \(\sigma\in\{\Car,\Cur,\Hom\}\), a primitive forcing shell
    \[
    \ForSpace{\sigma},
    \]
    a visible reduction map
    \[
    \Reduction{\sigma}:\ForSpace{\sigma}\to X_{\sigma},
    \]
    and shell-level current and equilibrium carriers
    \[
    \CurrentCarrier{\sigma}:\ForSpace{\sigma}\to Z_{\sigma}^{\mathrm{cur}},
    \qquad
    \EqCarrier{\sigma}:\ForSpace{\sigma}\to Z_{\sigma}^{\mathrm{eq}}.
    \]
    \item The product shell
    \[
    \PrimShell
    \equiv
    \ForSpace{\Car}\times \ForSpace{\Cur}\times \ForSpace{\Hom}.
    \]
    \item Fixed shell extractors
    \[
    \CurrentExtract{\sigma}:Z_{\sigma}^{\mathrm{cur}}\to Y_{\sigma}^{\mathrm{cur}},
    \qquad
    \EqExtract{\sigma}:Z_{\sigma}^{\mathrm{eq}}\to Y_{\sigma}^{\mathrm{eq}},
    \]
    together with the recorded shell composites
    \begin{equation}
    \PrimCurrent{\sigma}
    =
    \CurrentExtract{\sigma}\circ \CurrentCarrier{\sigma},
    \qquad
    \PrimTheta{\sigma}
    =
    \EqExtract{\sigma}\circ \EqCarrier{\sigma}.
    \label{eq:primcomposites}
    \end{equation}
    \item The recorded metric and ordinary-matter outputs
    \[
    \MetricOut:\PrimShell\to\{\text{observer metrics}\},
    \qquad
    \MatterOut:\PrimShell\times\Psi\to\{\text{observer stress tensors}\},
    \]
    satisfying
    \begin{equation}
    \MetricOut(\mathbf w)=\CompletedMetric,
    \qquad
    \MatterOut(\mathbf w,\psi)
    =
    T_{\mu\nu}^{\mathrm{matter}}[\CompletedMetric,\psi]
    \label{eq:fixedoutputs}
    \end{equation}
    for every \(\mathbf w\in\PrimShell\) and every matter configuration \(\psi\in\Psi\).
\end{enumerate}
\end{definition}

\begin{remark}
Definition \ref{def:fixed} is not reopened here. The primitive shell data and the benchmark one-metric / ordinary-matter outputs remain fixed continuation data. The work begins one layer deeper again: why the selector law that now carries the active line should itself be forced.
\end{remark}

\section{Target Selector Generator Law}

\begin{definition}[Zero-hidden selector generator law]
\label{def:generator}
Fix \(\sigma\in\{\Car,\Cur,\Hom\}\). A \emph{zero-hidden selector generator law} for sector \(\sigma\) consists of the following data.
\begin{enumerate}
    \item A finite-dimensional Euclidean selector-state space
    \[
    \GenSpace{\sigma},
    \]
    the observer-data space
    \[
    \DataSpace{\sigma}
    \equiv
    \VisibleAffine{\sigma}\times Z_{\sigma}^{\mathrm{cur}}\times Z_{\sigma}^{\mathrm{eq}},
    \]
    an ambient shell-response space
    \[
    \AmbientTarget{\sigma},
    \]
    and a compact convex selector-state body
    \[
    \GenBody{\sigma}\subset\GenSpace{\sigma}
    \]
    with nonempty interior.
    \item An affine observer-data readout
    \[
    \GenData{\sigma}:\GenSpace{\sigma}\to\DataSpace{\sigma}
    \]
    and an affine shell-response readout
    \[
    \GenShell{\sigma}:\GenSpace{\sigma}\to\AmbientTarget{\sigma}.
    \]
    \item A self-adjoint idempotent hidden selector
    \[
    \HiddenSel{\sigma}:\GenSpace{\sigma}\to\GenSpace{\sigma},
    \qquad
    {\HiddenSel{\sigma}}^{2}=\HiddenSel{\sigma},
    \qquad
    {\HiddenSel{\sigma}}^{*}=\HiddenSel{\sigma}.
    \]
    Define the selector-zero locus
    \[
    \ZeroHiddenSlice{\sigma}
    \equiv
    \ker\HiddenSel{\sigma},
    \]
    its translation space
    \[
    \ZeroHiddenVec{\sigma}
    \equiv
    \ker\HiddenSel{\sigma},
    \]
    and the hidden selector image
    \[
    \GapSpace{\sigma}
    \equiv
    \operatorname{im}\HiddenSel{\sigma}.
    \]
    \item A compact convex hidden body
    \[
    \GapBody{\sigma}\subset\GapSpace{\sigma}
    \]
    with
    \[
    0\in\operatorname{int}\GapBody{\sigma},
    \]
    such that the selector-state body splits exactly as
    \begin{equation}
    \GenBody{\sigma}
    =
    \left\{
    z+\eta:
    z\in\GenBody{\sigma}\cap\ZeroHiddenSlice{\sigma},
    \quad
    \eta\in\GapBody{\sigma}
    \right\}.
    \label{eq:gensplitting}
    \end{equation}
    \item A positive self-adjoint hidden-sector operator
    \[
    \HiddenOperator{\sigma}:\GapSpace{\sigma}\to\GapSpace{\sigma}
    \]
    satisfying
    \begin{equation}
    \ev{\eta,\HiddenOperator{\sigma}\eta}
    \ge
    \kappa_{\sigma}^{\mathrm{hid}}\,
    \|\eta\|^{2}
    \qquad
    \text{for all }\eta\in\GapSpace{\sigma}
    \label{eq:gengap}
    \end{equation}
    for some \(\kappa_{\sigma}^{\mathrm{hid}}>0\).
    \item A compact symmetry group
    \[
    \GenSym{\sigma}\subset
    O(\GenSpace{\sigma})\times
    O(\DataSpace{\sigma})\times
    O(\AmbientTarget{\sigma})
    \]
    and an orthogonal involution
    \[
    \GenTime{\sigma}\in
    O(\GenSpace{\sigma})\times
    O(\DataSpace{\sigma})\times
    O(\AmbientTarget{\sigma})
    \]
    preserving \(\GenBody{\sigma}\), \(\GenData{\sigma}\), \(\GenShell{\sigma}\), \(\HiddenSel{\sigma}\), \(\GapBody{\sigma}\), and \(\HiddenOperator{\sigma}\).
    \item A smooth embedding
    \[
    \jmath_{\sigma}:X_{\sigma}\hookrightarrow\VisibleAffine{\sigma}
    \]
    and a smooth zero-response shell realization
    \[
    \GenEmbed{\sigma}:\ForSpace{\sigma}\hookrightarrow\GenSpace{\sigma}
    \]
    whose image is exactly
    \[
    \GenBody{\sigma}\cap
    \ZeroHiddenSlice{\sigma}\cap
    \bigl(\GenShell{\sigma}\bigr)^{-1}(0),
    \]
    and such that
    \begin{equation}
    \GenData{\sigma}\bigl(\GenEmbed{\sigma}(w)\bigr)
    =
    \bigl(
    \jmath_{\sigma}(\Reduction{\sigma}(w)),
    \CurrentCarrier{\sigma}(w),
    \EqCarrier{\sigma}(w)
    \bigr)
    \label{eq:gencompat}
    \end{equation}
    for every \(w\in\ForSpace{\sigma}\).
    \item Fixed-data shell solvability on the selector-zero locus:
    \begin{equation}
    \GenData{\sigma}\bigl(
    \GenBody{\sigma}\cap\ZeroHiddenSlice{\sigma}
    \bigr)
    =
    \GenData{\sigma}\Bigl(
    \GenBody{\sigma}\cap
    \ZeroHiddenSlice{\sigma}\cap
    \bigl(\GenShell{\sigma}\bigr)^{-1}(0)
    \Bigr).
    \label{eq:gensolvability}
    \end{equation}
    \item The inert zero-hidden kernel
    \[
    \Kernel{\sigma}
    \equiv
    \ker\bigl(D\GenData{\sigma}\big|_{\ZeroHiddenVec{\sigma}}\bigr)
    \cap
    \ker\bigl(D\GenShell{\sigma}\big|_{\ZeroHiddenVec{\sigma}}\bigr)
    \]
    and the orthogonal projector
    \[
    \StateComp{\sigma}:
    \ZeroHiddenVec{\sigma}\to
    {\Kernel{\sigma}}^{\perp}\cap\ZeroHiddenVec{\sigma}.
    \]
    \item Selector-law coercivity on data-invisible directions: there exists \(\kappa_{\sigma}^{\mathrm{sel}}>0\) such that for every
    \[
    w\in\ForSpace{\sigma},
    \qquad
    \delta g\in T_{\GenEmbed{\sigma}(w)}\GenSpace{\sigma}
    \]
    satisfying
    \[
    D\GenData{\sigma}\,\delta g=0,
    \]
    one has
    \begin{equation}
    \|D\GenShell{\sigma}\,\delta g\|^{2}
    +
    \ev{
    \HiddenSel{\sigma}\delta g,
    \HiddenOperator{\sigma}\HiddenSel{\sigma}\delta g
    }
    \ge
    \kappa_{\sigma}^{\mathrm{sel}}
    \left(
    \|\StateComp{\sigma}( (I-\HiddenSel{\sigma})\delta g )\|^{2}
    +
    \|\HiddenSel{\sigma}\delta g\|^{2}
    \right).
    \label{eq:gencoercive}
    \end{equation}
\end{enumerate}
\end{definition}

\begin{remark}
Definition \ref{def:generator} records exactly the selector-layer objects that now require deeper justification. The present note does not enlarge that selector package. It asks whether those same objects can be forced by a smaller law beneath them.
\end{remark}

\section{Zero-Hidden Spectral Generator Law}

\begin{definition}[Zero-hidden spectral generator law]
\label{def:spectral}
Fix \(\sigma\in\{\Car,\Cur,\Hom\}\). A \emph{zero-hidden spectral generator law} for sector \(\sigma\) consists of the following data.
\begin{enumerate}
    \item A finite-dimensional Euclidean spectral state space
    \[
    \SpecSpace{\sigma},
    \]
    the observer-data space
    \[
    \DataSpace{\sigma}
    \equiv
    \VisibleAffine{\sigma}\times Z_{\sigma}^{\mathrm{cur}}\times Z_{\sigma}^{\mathrm{eq}},
    \]
    an ambient shell-response space
    \[
    \AmbientTarget{\sigma},
    \]
    and a compact convex spectral state body
    \[
    \SpecBody{\sigma}\subset\SpecSpace{\sigma}
    \]
    with nonempty interior.
    \item An affine observer-data readout
    \[
    \SpecData{\sigma}:\SpecSpace{\sigma}\to\DataSpace{\sigma}
    \]
    and an affine shell-response readout
    \[
    \SpecShell{\sigma}:\SpecSpace{\sigma}\to\AmbientTarget{\sigma}.
    \]
    \item A positive self-adjoint spectral generator
    \[
    \SpecOp{\sigma}:\SpecSpace{\sigma}\to\SpecSpace{\sigma},
    \qquad
    {\SpecOp{\sigma}}^{*}=\SpecOp{\sigma},
    \qquad
    \ev{u,\SpecOp{\sigma}u}\ge 0
    \]
    for all \(u\in\SpecSpace{\sigma}\). Define the zero-hidden spectral space
    \[
    \SpecZero{\sigma}\equiv\ker\SpecOp{\sigma},
    \]
    the hidden spectral space
    \[
    \SpecHidden{\sigma}
    \equiv
    \bigl(\SpecZero{\sigma}\bigr)^{\perp},
    \]
    and the orthogonal hidden spectral projector
    \[
    \SpecProj{\sigma}:\SpecSpace{\sigma}\to\SpecHidden{\sigma}.
    \]
    \item A compact convex hidden spectral body
    \[
    \SpecHiddenBody{\sigma}\subset\SpecHidden{\sigma}
    \]
    with
    \[
    0\in\operatorname{int}\SpecHiddenBody{\sigma},
    \]
    such that the spectral state body splits exactly as
    \begin{equation}
    \SpecBody{\sigma}
    =
    \left\{
    z+\eta:
    z\in\SpecBody{\sigma}\cap\SpecZero{\sigma},
    \quad
    \eta\in\SpecHiddenBody{\sigma}
    \right\}.
    \label{eq:specsplitting}
    \end{equation}
    \item An exact hidden-band gap: there exists \(\lambda_{\sigma}^{\mathrm{sp}}>0\) such that
    \begin{equation}
    \ev{\eta,\SpecOp{\sigma}\eta}
    \ge
    \lambda_{\sigma}^{\mathrm{sp}}\,
    \|\eta\|^{2}
    \qquad
    \text{for all }\eta\in\SpecHidden{\sigma}.
    \label{eq:specgap}
    \end{equation}
    \item A compact symmetry group
    \[
    \SpecSym{\sigma}\subset
    O(\SpecSpace{\sigma})\times
    O(\DataSpace{\sigma})\times
    O(\AmbientTarget{\sigma})
    \]
    and an orthogonal involution
    \[
    \SpecTime{\sigma}\in
    O(\SpecSpace{\sigma})\times
    O(\DataSpace{\sigma})\times
    O(\AmbientTarget{\sigma})
    \]
    preserving \(\SpecBody{\sigma}\), \(\SpecData{\sigma}\), \(\SpecShell{\sigma}\), and \(\SpecOp{\sigma}\).
    \item A smooth embedding
    \[
    \jmath_{\sigma}:X_{\sigma}\hookrightarrow\VisibleAffine{\sigma}
    \]
    and a smooth zero-response shell realization
    \[
    \SpecEmbed{\sigma}:\ForSpace{\sigma}\hookrightarrow\SpecSpace{\sigma}
    \]
    whose image is exactly
    \[
    \SpecBody{\sigma}\cap
    \SpecZero{\sigma}\cap
    \bigl(\SpecShell{\sigma}\bigr)^{-1}(0),
    \]
    and such that
    \begin{equation}
    \SpecData{\sigma}\bigl(\SpecEmbed{\sigma}(w)\bigr)
    =
    \bigl(
    \jmath_{\sigma}(\Reduction{\sigma}(w)),
    \CurrentCarrier{\sigma}(w),
    \EqCarrier{\sigma}(w)
    \bigr)
    \label{eq:speccompat}
    \end{equation}
    for every \(w\in\ForSpace{\sigma}\).
    \item Fixed-data shell solvability on the zero-hidden spectral space:
    \begin{equation}
    \SpecData{\sigma}\bigl(
    \SpecBody{\sigma}\cap\SpecZero{\sigma}
    \bigr)
    =
    \SpecData{\sigma}\Bigl(
    \SpecBody{\sigma}\cap
    \SpecZero{\sigma}\cap
    \bigl(\SpecShell{\sigma}\bigr)^{-1}(0)
    \Bigr).
    \label{eq:specsolvability}
    \end{equation}
    \item The inert zero-hidden kernel
    \[
    \Kernel{\sigma}
    \equiv
    \ker\bigl(D\SpecData{\sigma}\big|_{\SpecZero{\sigma}}\bigr)
    \cap
    \ker\bigl(D\SpecShell{\sigma}\big|_{\SpecZero{\sigma}}\bigr)
    \]
    and the orthogonal projector
    \[
    \StateComp{\sigma}:
    \SpecZero{\sigma}\to
    {\Kernel{\sigma}}^{\perp}\cap\SpecZero{\sigma}.
    \]
    \item Spectral coercivity on data-invisible directions: there exists \(\kappa_{\sigma}^{\mathrm{sp}}>0\) such that for every
    \[
    w\in\ForSpace{\sigma},
    \qquad
    \delta u\in T_{\SpecEmbed{\sigma}(w)}\SpecSpace{\sigma}=\SpecSpace{\sigma}
    \]
    satisfying
    \[
    D\SpecData{\sigma}\,\delta u=0,
    \]
    one has
    \begin{equation}
    \|D\SpecShell{\sigma}\,\delta u\|^{2}
    +
    \ev{\delta u,\SpecOp{\sigma}\delta u}
    \ge
    \kappa_{\sigma}^{\mathrm{sp}}
    \left(
    \|\StateComp{\sigma}\bigl((I-\SpecProj{\sigma})\delta u\bigr)\|^{2}
    +
    \|\SpecProj{\sigma}\delta u\|^{2}
    \right).
    \label{eq:speccoercive}
    \end{equation}
\end{enumerate}
\end{definition}

\begin{remark}
Definition \ref{def:spectral} is strictly smaller than Definition \ref{def:generator} in the precise way now needed. The hidden selector and the hidden-sector operator are no longer stipulated separately. They are both extracted from one self-adjoint spectral generator with an exact zero band and a positively gapped hidden band.
\end{remark}

\section{Canonical Selector Law Induced by the Spectral Generator}

\begin{definition}[Canonical selector law induced by the spectral generator]
\label{def:canonical}
Assume Definition \ref{def:spectral}. Define the canonical zero-hidden selector generator law by:
\begin{enumerate}
    \item
    \[
    \GenSpace{\sigma}\equiv\SpecSpace{\sigma},
    \qquad
    \GenBody{\sigma}\equiv\SpecBody{\sigma},
    \qquad
    \GenData{\sigma}\equiv\SpecData{\sigma},
    \qquad
    \GenShell{\sigma}\equiv\SpecShell{\sigma}.
    \]
    \item The hidden selector is the orthogonal spectral projector
    \begin{equation}
    \HiddenSel{\sigma}\equiv\SpecProj{\sigma}.
    \label{eq:selectorproj}
    \end{equation}
    \item
    \[
    \ZeroHiddenSlice{\sigma}\equiv\SpecZero{\sigma},
    \qquad
    \ZeroHiddenVec{\sigma}\equiv\SpecZero{\sigma},
    \qquad
    \GapSpace{\sigma}\equiv\SpecHidden{\sigma},
    \qquad
    \GapBody{\sigma}\equiv\SpecHiddenBody{\sigma}.
    \]
    \item The hidden-sector operator is the restriction
    \begin{equation}
    \HiddenOperator{\sigma}
    \equiv
    \SpecOp{\sigma}\big|_{\SpecHidden{\sigma}}.
    \label{eq:hiddenrestriction}
    \end{equation}
    \item
    \[
    \GenSym{\sigma}\equiv\SpecSym{\sigma},
    \qquad
    \GenTime{\sigma}\equiv\SpecTime{\sigma},
    \qquad
    \GenEmbed{\sigma}\equiv\SpecEmbed{\sigma}.
    \]
\end{enumerate}
\end{definition}

\begin{lemma}[The spectral generator determines the hidden selector]
\label{lem:projection}
Assume Definition \ref{def:spectral}. Then
\[
\SpecSpace{\sigma}
=
\SpecZero{\sigma}\oplus\SpecHidden{\sigma},
\qquad
\SpecHidden{\sigma}=\operatorname{im}\SpecOp{\sigma},
\]
orthogonally, and the projector \(\SpecProj{\sigma}\) is self-adjoint and idempotent with
\[
\ker\SpecProj{\sigma}=\SpecZero{\sigma},
\qquad
\operatorname{im}\SpecProj{\sigma}=\SpecHidden{\sigma}.
\]
\end{lemma}

\begin{proof}
Because \(\SpecOp{\sigma}\) is self-adjoint on a finite-dimensional Euclidean space, one has the orthogonal decomposition
\[
\SpecSpace{\sigma}
=
\ker\SpecOp{\sigma}\oplus \bigl(\ker\SpecOp{\sigma}\bigr)^{\perp}.
\]
Thus \(\SpecSpace{\sigma}=\SpecZero{\sigma}\oplus\SpecHidden{\sigma}\) orthogonally. For any self-adjoint operator on a finite-dimensional inner-product space,
\[
\operatorname{im}\SpecOp{\sigma}
=
\bigl(\ker\SpecOp{\sigma}\bigr)^{\perp}
=
\SpecHidden{\sigma}.
\]
The orthogonal projector onto \(\SpecHidden{\sigma}\) is therefore self-adjoint and idempotent, with kernel \(\SpecZero{\sigma}\) and image \(\SpecHidden{\sigma}\).
\end{proof}

\begin{proposition}[The spectral law induces a zero-hidden selector generator law]
\label{prop:generator}
Assume Definitions \ref{def:spectral} and \ref{def:canonical}. Then the canonical data of Definition \ref{def:canonical} satisfy every requirement of Definition \ref{def:generator}.
\end{proposition}

\begin{proof}
Items (1) and (2) of Definition \ref{def:generator} are immediate from Definition \ref{def:canonical}. By Lemma \ref{lem:projection}, the operator \(\HiddenSel{\sigma}=\SpecProj{\sigma}\) is self-adjoint and idempotent, with
\[
\ZeroHiddenSlice{\sigma}
=
\ker\HiddenSel{\sigma}
=
\SpecZero{\sigma},
\qquad
\GapSpace{\sigma}
=
\operatorname{im}\HiddenSel{\sigma}
=
\SpecHidden{\sigma}.
\]
This gives item (3).

The compact convex hidden body is \(\GapBody{\sigma}=\SpecHiddenBody{\sigma}\), and the selector-state body splitting \eqref{eq:gensplitting} is exactly the spectral body splitting \eqref{eq:specsplitting}. This gives item (4). For item (5), \(\HiddenOperator{\sigma}\) is the restriction of the positive self-adjoint operator \(\SpecOp{\sigma}\) to \(\SpecHidden{\sigma}\), so it is positive and self-adjoint, and \eqref{eq:gengap} follows from \eqref{eq:specgap}.

The symmetry and time-reversal requirements follow because every element of \(\SpecSym{\sigma}\) and \(\SpecTime{\sigma}\) preserves \(\SpecOp{\sigma}\), hence preserves both \(\SpecZero{\sigma}\) and \(\SpecHidden{\sigma}\), and therefore also the orthogonal projector \(\SpecProj{\sigma}\), the hidden body \(\SpecHiddenBody{\sigma}\), and the restricted operator \(\HiddenOperator{\sigma}\). Thus item (6) holds.

The shell realization and data compatibility in item (7) are exactly \eqref{eq:speccompat}, because
\[
\GenBody{\sigma}\cap\ZeroHiddenSlice{\sigma}\cap\bigl(\GenShell{\sigma}\bigr)^{-1}(0)
=
\SpecBody{\sigma}\cap\SpecZero{\sigma}\cap\bigl(\SpecShell{\sigma}\bigr)^{-1}(0).
\]
Fixed-data shell solvability \eqref{eq:gensolvability} is exactly \eqref{eq:specsolvability}. The inert zero-hidden kernel and orthogonal projector in item (9) are defined by the same formulas, since
\[
\GenData{\sigma}=\SpecData{\sigma},
\qquad
\GenShell{\sigma}=\SpecShell{\sigma},
\qquad
\ZeroHiddenVec{\sigma}=\SpecZero{\sigma}.
\]

For item (10), let
\[
w\in\ForSpace{\sigma},
\qquad
\delta g\in T_{\GenEmbed{\sigma}(w)}\GenSpace{\sigma}=\SpecSpace{\sigma}
\]
with \(D\GenData{\sigma}\,\delta g=0\). Since \(\HiddenSel{\sigma}=\SpecProj{\sigma}\), one has
\[
(I-\HiddenSel{\sigma})\delta g=(I-\SpecProj{\sigma})\delta g\in\SpecZero{\sigma},
\qquad
\HiddenSel{\sigma}\delta g=\SpecProj{\sigma}\delta g\in\SpecHidden{\sigma}.
\]
Because \(\SpecOp{\sigma}\) vanishes on \(\SpecZero{\sigma}\) and agrees with \(\HiddenOperator{\sigma}\) on \(\SpecHidden{\sigma}\),
\[
\ev{\delta g,\SpecOp{\sigma}\delta g}
=
\ev{\SpecProj{\sigma}\delta g,\SpecOp{\sigma}\SpecProj{\sigma}\delta g}
=
\ev{\HiddenSel{\sigma}\delta g,\HiddenOperator{\sigma}\HiddenSel{\sigma}\delta g}.
\]
Applying \eqref{eq:speccoercive} therefore gives \eqref{eq:gencoercive} with \(\kappa_{\sigma}^{\mathrm{sel}}=\kappa_{\sigma}^{\mathrm{sp}}\). All requirements of Definition \ref{def:generator} are satisfied.
\end{proof}

\section{Forced Constituents and Uniqueness}

\begin{proposition}[All selector constituents are forced by the spectral law]
\label{prop:forced}
Assume Definition \ref{def:spectral}. Then the following selector-level constituents are fixed by the spectral law itself.
\begin{enumerate}
    \item The selector-state body is necessarily the spectral state body \(\GenBody{\sigma}=\SpecBody{\sigma}\).
    \item The affine observer-data and shell-response readouts are necessarily \(\GenData{\sigma}=\SpecData{\sigma}\) and \(\GenShell{\sigma}=\SpecShell{\sigma}\).
    \item The hidden selector is necessarily the orthogonal projector onto the hidden spectral band:
    \[
    \HiddenSel{\sigma}=\SpecProj{\sigma}.
    \]
    \item The selector-zero locus and translation space are necessarily
    \[
    \ZeroHiddenSlice{\sigma}=\ZeroHiddenVec{\sigma}=\ker\SpecOp{\sigma},
    \]
    while the hidden selector image is necessarily
    \[
    \GapSpace{\sigma}=\operatorname{im}\SpecOp{\sigma}.
    \]
    \item The selector-zero / hidden body splitting is necessarily the spectral splitting \eqref{eq:specsplitting}, the hidden body is necessarily \(\SpecHiddenBody{\sigma}\), and the positive hidden-sector operator is necessarily the restriction \(\SpecOp{\sigma}\big|_{\operatorname{im}\SpecOp{\sigma}}\).
    \item The symmetry and time-reversal structure, the exact zero-response shell realization, fixed-data shell solvability on the selector-zero locus, the inert zero-hidden kernel and projector, and the selector-law coercivity inequality are all induced by the same spectral law and are not additional selector freedom.
\end{enumerate}
\end{proposition}

\begin{proof}
Items (1) and (2) are immediate from Definition \ref{def:canonical}: once the spectral state body and the two affine readouts are fixed, no second selector-state body or second pair of selector readouts is compatible with the same spectral law. Item (3) is Lemma \ref{lem:projection}. Item (4) follows from the same lemma together with the self-adjoint finite-dimensional identity
\[
\operatorname{im}\SpecOp{\sigma}
=
\bigl(\ker\SpecOp{\sigma}\bigr)^{\perp}.
\]
Item (5) is exactly the content of Definition \ref{def:canonical}: the body splitting, hidden body, and hidden operator are the corresponding spectral objects and carry no further freedom. Item (6) is the content of Proposition \ref{prop:generator}: every remaining selector constituent descends directly from the same spectral package.
\end{proof}

\begin{theorem}[Zero-hidden selector generator law necessity / uniqueness from spectral law]
\label{thm:necessity}
Assume Definition \ref{def:spectral}. Then:
\begin{enumerate}
    \item the canonical construction of Definition \ref{def:canonical} yields a zero-hidden selector generator law in the sense of Definition \ref{def:generator};
    \item every selector-level object singled out by the previous ambient-dynamics theorem is forced by the spectral law: the compact convex selector-state body, the affine observer-data and shell-response readouts, the self-adjoint idempotent hidden selector, the selector-zero locus and hidden selector image, the selector-zero / hidden body splitting, the hidden factor and positive hidden-sector operator, the symmetry and time-reversal structure, the exact zero-response shell realization on the selector-zero locus, fixed-data shell solvability on that locus, the inert zero-hidden kernel and projector, and the selector-law coercivity inequality;
    \item any other selector package compatible with the same spectral law differs from the canonical one only by the obvious orthogonal identification that preserves the zero/hidden spectral splitting and the recorded data and shell coordinates;
    \item consequently the selector-law layer is not an independent remaining freedom once the spectral law is fixed.
\end{enumerate}
\end{theorem}

\begin{proof}
Item (1) is Proposition \ref{prop:generator}. Item (2) is Proposition \ref{prop:forced}. For item (3), any compatible reparameterization must preserve the orthogonal decomposition
\[
\SpecSpace{\sigma}
=
\ker\SpecOp{\sigma}\oplus\operatorname{im}\SpecOp{\sigma},
\]
must preserve the spectral generator on each factor, and must leave the recorded data and shell coordinates unchanged. That leaves only the obvious orthogonal identification respecting the zero/hidden spectral splitting. Item (4) is the routing consequence of the previous items.
\end{proof}

\section{Benchmark Preservation}

\begin{theorem}[The derived selector law preserves the same one-metric benchmark package]
\label{thm:outputs}
Assume Definition \ref{def:spectral}. Then the benchmark output statements of Definition \ref{def:fixed} remain unchanged:
\[
\MetricOut(\mathbf w)=\CompletedMetric,
\qquad
\MatterOut(\mathbf w,\psi)
=
T_{\mu\nu}^{\mathrm{matter}}[\CompletedMetric,\psi]
\]
for every \(\mathbf w\in\PrimShell\) and every matter configuration \(\psi\in\Psi\). Consequently the deeper spectral law and the induced selector law introduce neither a second operative metric nor an enlarged low-energy matter law.
\end{theorem}

\begin{proof}
The present note changes only the upstream origin of the selector data. It does not alter the primitive shell \(\PrimShell\), the recorded metric map \(\MetricOut\), or the recorded matter map \(\MatterOut\). Those remain the fixed continuation data of Definition \ref{def:fixed}, so \eqref{eq:fixedoutputs} continues to hold exactly.
\end{proof}

\section{Main Theorem}

\begin{theorem}[The remaining burden moves below the selector-law layer]
\label{thm:main}
Assume the fixed primitive continuation data of Definition \ref{def:fixed} and, for each sector \(\sigma\in\{\Car,\Cur,\Hom\}\), a zero-hidden spectral generator law in the sense of Definition \ref{def:spectral}. Then:
\begin{enumerate}
    \item the zero-hidden selector generator law of Definition \ref{def:generator} is canonically induced by the spectral law;
    \item the compact convex selector-state body, the affine observer-data and shell-response readouts, the self-adjoint idempotent hidden selector, the selector-zero locus and hidden selector image, the selector-zero / hidden body splitting, the hidden factor and positive hidden-sector operator, the symmetry and time-reversal structure, the exact zero-response shell realization on the selector-zero locus, fixed-data shell solvability on that locus, the inert zero-hidden kernel and projector, and the selector-law coercivity inequality are all forced by that same spectral law;
    \item the one operative metric remains exactly \(\CompletedMetric\), the ordinary matter route remains exactly the standard one through that metric, there is no second operative metric, and the low-energy matter law is not enlarged;
    \item consequently the benchmark-facing burden after the previous ambient-dynamics theorem is discharged in the stronger form now available on the active line: the zero-hidden selector generator law is no longer an unexplained stopping point, but a necessary reduction of a smaller zero-hidden spectral generator law.
\end{enumerate}
\end{theorem}

\begin{proof}
Item (1) is Theorem \ref{thm:necessity}(1). Item (2) is Theorem \ref{thm:necessity}(2). Item (3) is Theorem \ref{thm:outputs}. Item (4) is the combined routing consequence.
\end{proof}

\begin{corollary}[What remains open after this theorem]
\label{cor:next}
After Theorem \ref{thm:main}, the remaining honest burden is no longer to justify the zero-hidden selector generator law itself on the active primitive-observer-reduction line. The next honest burden is deeper again: derive or justify the zero-hidden spectral generator law from still more primitive substrate dynamics, or else prove a sharper inevitability theorem that removes even that remaining spectral-generator-level freedom while preserving the same benchmark one-metric package and claim boundary.
\end{corollary}

\begin{proof}
Theorem \ref{thm:main} converts the selector-law layer from a remaining benchmark-facing burden into a canonical reduction of a deeper spectral law. What remains open is why that spectral law itself should hold, rather than becoming the new disciplined stopping point.
\end{proof}

\section{Conclusion}

The positive result of this note is deliberately narrower than a completed final microphysical derivation and stronger than another deeper same-output relay. The previous theorem had already shown that the ambient-balance layer is a forced reduction of a selector law. The present note goes one honest step further by showing that the selector layer itself is not left hanging either. It is canonically induced by a smaller zero-hidden spectral generator law.

That gain directly addresses the precise objects that had become the next burden: the compact convex selector-state body, the affine observer-data and shell-response readouts, the self-adjoint idempotent hidden selector, the selector-zero locus and hidden selector image, the selector-zero / hidden body splitting, the hidden factor and positive hidden-sector operator, the symmetry and time-reversal structure, the exact zero-response shell realization on the selector-zero locus, fixed-data shell solvability on that locus, the inert zero-hidden kernel and projector, and the selector-law coercivity inequality. Within the spectral-law class those are no longer optional inserted ingredients. They are forced outputs, unique up to the harmless orthogonal identification preserving the zero/hidden spectral splitting.

The benchmark boundary remains fixed throughout. The same primitive shell remains the shell carrying the observer outputs, so the one operative metric is still exactly \(\CompletedMetric\), the ordinary matter route is unchanged, there is no second operative metric, and the low-energy matter law is not enlarged. This is therefore a disciplined upstream derivational gain, not a final completion claim. The note still does not derive the spectral law from ultimate substrate microphysics. That is the next honest open burden. But the selector-law layer is no longer the stopping point on the active line.

\end{document}
